\subsection{Three regular six-element monoids and their full homotopy types}

\begin{theorem}[A group above a rectangular band]\label{E:group-band-cofiber}
Let $G$ be a finite group, $I$ a nonempty finite left $G$-set, and $J$ a
nonempty finite right $G$-set. On
\[
 M=G\amalg(I\times J)
\]
retain the multiplication of $G$ and define
\[
 (i,j)(i',j')=(i,j'),\qquad
 g(i,j)=(gi,j),\qquad (i,j)g=(i,jg).
\]
This is a finite regular monoid. Give $I*J$ the left $G$-action induced by
$g i=gi$ and $g j=jg^{-1}$. There is a homotopy equivalence
\[
 BM\simeq\operatorname{hocofib}
 \bigl(EG\times_G(I*J)\longrightarrow BG\bigr),
\]
where the displayed map is the bundle projection. In particular this is a
homotopy-type assertion in every degree, not an inference from initial
homology groups.
\end{theorem}

\begin{proof}
The two actions commute on $I\times J$, and the rectangular-band product
uses only the outer coordinates. Checking the possible placements of group
and band factors therefore proves associativity. The identity of $G$ is the
identity of $M$; every band element is idempotent and every group element is
invertible, so $M$ is regular.

\textbf{Step 1: a category with three objects.}
Let $\mathcal C$ have objects $A,B,C$, endomorphism groups $1,G,1$,
respectively, and nonempty cross hom-sets
\[
 \mathcal C(A,B)=I,\qquad \mathcal C(B,C)=J,\qquad
 \mathcal C(A,C)=\{t\}.
\]
Composition with $G$ is given by the specified actions, and every composite
$j\circ i$ is $t$. All other cross hom-sets are empty. These rules are
associative, including $j\circ g\circ i$, since both associations of this
last expression give $t$.

Choose $i_0\in I$, $j_0\in J$, and put $e=(i_0,j_0)\in M$.
The full subcategory $\mathcal D$ of the Karoubi category of $M$ on $1,e$
is equivalent to the entire Karoubi category: its only other objects are
the band idempotents, each isomorphic to $e$. The standard equivalence
$B\operatorname{Kar}(M)\simeq BM$ was justified in
Lemma~\ref{A:minimal-group}. There is a functor
\[
 F:\mathcal C\longrightarrow\mathcal D,
 \qquad F(A)=F(C)=e,\quad F(B)=1,
\]
given on arrows by
\[
 F(g)=g,\qquad F(i)=(i,j_0),\qquad
 F(j)=(i_0,j),\qquad F(t)=e.
\]
For example $F(j)F(i)=(i_0,j_0)=e$, so these are exactly the required
composition rules.

\textbf{Step 2: the two comma categories are contractible.}
We check $d\mathbin{\downarrow}F$ for both objects $d$ of $\mathcal D$,
using the convention that its objects are $(X,a:d\to F(X))$.
For $d=e$ the object $(A,e)$ is initial: the objects over $B$ are indexed
by $(i,j_0)$, and the unique arrow from $(A,e)$ to that object is $i$.
There is also the unique arrow $t$ to the object over $C$.

For $d=1$ the objects can be written
\[
 A_j,\quad B_g,\quad B_{i,j},\quad C_j
 \qquad (g\in G,\ i\in I,\ j\in J).
\]
Let $\mathcal Q_0$ be the full subcategory on all $B_g$ and $C_j$.
It has initial object $B_1$: the arrow from $B_1$ to $B_g$ is $g$,
and the arrow to $C_j$ is $j$. For each $j\in J$, let
$\mathcal Q_j$ be the full subcategory on
$A_j$, all $B_{i,j}$, and $C_j$. It has initial object $A_j$.
Indeed the arrow $A_j\to B_{i,j}$ is the unique arrow $i$, and
$A_j\to C_j$ is $t$.

These subcategories contain every arrow of $1\mathbin{\downarrow}F$.
Their only intersections are
$\mathcal Q_0\cap\mathcal Q_j=\{C_j\}$, and each $C_j$ is a sink.
Thus every composable string lies in one of the displayed subcategories.
The nerve is the union of their nerves, with each contractible
$B\mathcal Q_j$ attached to the contractible $B\mathcal Q_0$ at the
single vertex $C_j$. All these are CW subcomplexes, so the union is
contractible. Quillen's Theorem A, in its comma-category form
\cite[Theorem A]{quillen1973}, now gives
\[
 B\mathcal C\simeq B\mathcal D\simeq BM.
\]

\textbf{Step 3: the cofiber model, with its gluing specified.}
Let $L\subset B\mathcal C$ be the closed edge corresponding to $t:A\to C$.
It is a contractible CW subcomplex, so collapsing it is a homotopy
equivalence. A simplex which visits $B$ consists of an initial arrow
$i\in I$, if it starts at $A$, a string
\[
 g_1,\ldots,g_r\in G
\]
at $B$, and a final arrow $j\in J$, if it ends at $C$.
The strings visiting both $A$ and $C$ are consequently indexed by
$I\times G^r\times J$. Their faces inside the $B$ block multiply adjacent
group elements, while the first and last faces of that block replace
$i$ by $g_1i$ and $j$ by $jg_r$, respectively. This is the two-sided bar
description of $(I\times J)_{hG}$ for the diagonal action
$g(i,j)=(gi,jg^{-1})$.

For clarity, an explicit identification with the action-groupoid nerve
sends $(i,g_1,\ldots,g_r,j)$ to the group string
$(g_1,\ldots,g_r)$ starting at
\[
 (i,jP),\qquad P=g_r\cdots g_1.
\]
Its endpoint is $(Pi,j)$. The first and last face formulas above now agree
exactly with those of the action-groupoid nerve. Forgetting $j$ or $i$
gives the corresponding bar models of $I_{hG}$ or $J_{hG}$, and forgetting
both gives $BG$.

The join $I*J$ is the double mapping cylinder of
$I\leftarrow I\times J\to J$. The Borel construction preserves this
gluing, so its total space is the double mapping cylinder
\[
 I_{hG}\ \longleftarrow\ (I\times J)_{hG}\
 \longrightarrow\ J_{hG}
\]
over $BG$. The preceding simplex description identifies
$B\mathcal C/L$ with the mapping cone of its projection to $BG$.
One can check the identification directly in barycentric coordinates.
For a simplex visiting both endpoints, write its endpoint weights as
$t_A,t_C$ and its $B$-block weights as $s_0,\ldots,s_r$.
The cone parameter is $v=t_A+t_C$; for $0<v<1$ the join parameter is
$t_C/v$ and the bar-simplex weights are $s_k/(1-v)$.
At $v=0$ this is the projection to $BG$; at $v=1$ the simplex lies in
the collapsed edge $L$. At either end of the join parameter the unused
$i$ or $j$ is forgotten. The internal faces have exactly the bar gluing
just checked. The formulas therefore extend across all faces and
degeneracies and give the asserted mapping-cone model. Combining this
with Steps 1 and 2 proves the theorem.
\end{proof}

\begin{corollary}[The three two-point actions]\label{E:three-Thom}
Let $G=C_2$ and let each of $I,J$ have two elements. Up to interchanging
left and right, the corresponding monoids in
Theorem~\ref{E:group-band-cofiber} have the following homotopy types:
\[
\begin{array}{c|c|c|c}
\text{action on }I&\text{action on }J&\text{bundle over }\mathbb RP^\infty
 &BM\\\hline
\text{trivial}&\text{trivial}&\varepsilon\oplus\varepsilon
 &S^2\vee\Sigma^2\mathbb RP^\infty\\
\text{trivial}&\text{interchange}&\varepsilon\oplus\gamma
 &\Sigma\mathbb RP^\infty\\
\text{interchange}&\text{interchange}&\gamma\oplus\gamma
 &\mathbb RP^\infty/\mathbb RP^1
\end{array}
\]
Here $\gamma$ is the tautological real line bundle and $\varepsilon$ the
trivial real line bundle. These are, respectively, the supplied
six-element classes $138,152,190$; changing the single nontrivial action
from $J$ to $I$ gives the opposite of class $152$ and the same homotopy type.
All three monoids have trivial universal group.
\end{corollary}

\begin{proof}
A two-point $C_2$-set is the unit sphere of the trivial or sign
one-dimensional real representation. If these representations are $V,W$,
the usual join coordinates give an equivariant homeomorphism
\[
 S(V)*S(W)\cong S(V\oplus W).
\]
Consequently the projection in Theorem~\ref{E:group-band-cofiber} is
the unit-sphere bundle of the associated two-plane bundle $\xi$ over
$BC_2=\mathbb RP^\infty$. The disk bundle retracts to the base, so the
homotopy cofiber of its sphere-bundle projection is the Thom space
$\operatorname{Th}(\xi)$.

For the trivial two-plane bundle this is
$S^2\wedge(\mathbb RP^\infty)_+\simeq
S^2\vee\Sigma^2\mathbb RP^\infty$. For the second bundle,
\[
 \operatorname{Th}(\varepsilon\oplus\gamma)
 \cong\Sigma\operatorname{Th}(\gamma)
 \cong\Sigma\mathbb RP^\infty.
\]
For the third it is $\mathbb RP^\infty/\mathbb RP^1$.
The last two identifications use the explicit finite-dimensional
Thom-space formula
\[
 \operatorname{Th}(k\gamma|_{\mathbb RP^n})
 \cong\mathbb RP^{n+k}/\mathbb RP^{k-1},
\]
followed by the compatible direct limit. To see this formula, regard
$\mathbb RP^{n+k}\setminus\mathbb RP^{k-1}$ as the total space of
$k\gamma$ by splitting coordinates into a nonzero vector in
$\mathbb R^{n+1}$ and a vector in $\mathbb R^k$; simultaneous sign change
is exactly the transition rule for $k\gamma$. Its one-point
compactification is the indicated quotient.

The labels can be checked without relying on a census: with band elements
ordered as $(0,0),(1,0),(0,1),(1,1)$, multiplication is exactly the table
of class $138$ in the first case and of class $190$ in the third.
In the middle case interchange the second and third band labels to obtain
the displayed table of class $152$. Opposite categories have homeomorphic
realizations by reversing simplex coordinates, which handles the reversed
single-action case. Finally every band element is idempotent and hence
trivial in the universal group; $g(i,j)=(gi,j)$ then kills the generator
$g\in C_2$, including when its action is trivial.
\end{proof}

\begin{remark}[Finite certificates and their scope]\label{E:small-infinite-certificate}
The verification code is included as \texttt{a\_small\_infinite.py}.
It checks all four pairs of two-point actions. It verifies the monoid tables, every defined
composition and associativity equation of the three-object comparison
category and its two-object Karoubi category, the comparison functor,
both comma categories, their displayed initial objects and their
point-intersection decomposition. It also records exact table
isomorphisms with classes $138,152,190$ or their opposites.
The output is \texttt{code/outputs/a\_small\_infinite.json}.
These finite checks certify the category data used in the proof;
the all-degree homotopy identifications follow from the preceding
arguments, not from finite truncations of the nerve.
\end{remark}

\subsection{The complete answer for monoids of order at most six}
\begin{theorem}\label{E:all-six-types}
Every monoid of order at most six has exactly one of the homotopy types
in the following table. The entries count isomorphism-or-opposite classes
of monoids of the indicated order, not homotopy types or labelled tables.
Thus every small monoid with unbounded homology is identified in all
degrees, as requested in E5.
\[
\begin{array}{l|rrrrrr}
\text{homotopy type}&1&2&3&4&5&6\\\hline
* & 1 & 1 & 4 & 19 & 126 & 1159 \\
S^2 & 0 & 0 & 0 & 0 & 1 & 2 \\
BC_2 & 0 & 1 & 1 & 5 & 22 & 172 \\
BC_3 & 0 & 0 & 1 & 1 & 4 & 22 \\
BC_4 & 0 & 0 & 0 & 1 & 1 & 6 \\
B(C_2\times C_2) & 0 & 0 & 0 & 1 & 1 & 5 \\
BC_5 & 0 & 0 & 0 & 0 & 1 & 1 \\
BC_6 & 0 & 0 & 0 & 0 & 0 & 1 \\
BS_3 & 0 & 0 & 0 & 0 & 0 & 1 \\
S^2\vee\Sigma^2\RP^\infty & 0 & 0 & 0 & 0 & 0 & 1 \\
\Sigma\RP^\infty & 0 & 0 & 0 & 0 & 0 & 1 \\
\RP^\infty/\RP^1 & 0 & 0 & 0 & 0 & 0 & 1 \\
\bigvee_{n\geq2}S^n & 0 & 0 & 0 & 0 & 0 & 1 \\
\hline
\text{total}&1&2&6&27&156&1373
\end{array}
\]
\end{theorem}
\begin{proof}
The independent exhaustive tables of Section~\ref{sec:C} are reduced
recursively whenever an idempotent $e$ satisfies $eMe=Me$ or $eMe=eM$.
The explicit Karoubi comparison in Section~\ref{sec:A-census} proves
that each such reduction preserves the entire homotopy type: that
comparison uses only the stated one-sided equality, and applies to
the corner monoid even when it is not a group. The
terminal cores are the groups shown, the trivial monoid, the five-element
rectangular-band sphere monoid, the three group-band monoids just
identified, the Sweeney monoid, and one further six-element band.
The program \texttt{e\_small\_census.py} assigns every one of the $1565$
census classes to its core, checks the exact table isomorphisms for the
exceptional cores, and records each assignment in its JSON output.
There is no use of initial Betti numbers to infer an infinite homotopy type.

For clarity, the last six-element band has labels $1,a,c,b,d,q$, with
$\{a,c,b,d\}$ the $2\times2$ band, $qm=am$ and $mq=ma$ for band elements,
and $q^2=q$. A complete resolution of the trivial left $R=\Z M$ module is
\[
 0\longrightarrow Ra\xrightarrow{\cdot(b-a)}R(1-q)
 \xrightarrow{\cdot(c-a)}Ra\longrightarrow\Z\longrightarrow0.
\]
The respective integral bases are $(a,c)$,
$(1-q,b-a,d-c)$, and $(a,c)$. The rightmost differential sends $1-q$
to $c-a$ and the two other basis vectors to zero; the leftmost sends
$a,c$ to $b-a,d-c$. These formulas prove exactness integrally.
Tensoring with the right trivial module gives $\Z,0,\Z$ in degrees
$0,1,2$. Every element is idempotent, so the universal group is trivial.
The Hurewicz generator $S^2\to BM$ is therefore a homology equivalence
between simply connected CW spaces, hence a homotopy equivalence.
The independent certificate is
\texttt{b\_resolution\_root\_S2\_order6.json}.
The five-element core is also covered by Theorem~\ref{E:group-band-cofiber}
with trivial $G$ and two-point $I,J$, giving the suspension of a circle.
All terminal cores have now been identified by all-degree proofs.
Nontrivial finite groups, the three displayed Thom spaces, and the
infinite wedge have unbounded homology; the point and $S^2$ do not.
\end{proof}

\begin{corollary}\label{E:six-rings}
The three nongroup regular infinite types have zero positive cup products
in the first two cases $S^2\vee\Sigma^2\RP^\infty$ and
$\Sigma\RP^\infty$. In the third case
\[
 H^*(\RP^\infty/\RP^1;\F_2)
 \cong\F_2\oplus t^2\F_2[t],\qquad |t|=1.
\]
Its integral degree-two generator has every power nonzero; powers of
exponent at least two have order two.
\end{corollary}
\begin{proof}
Suspensions have zero products of positive reduced classes, and wedge
summands have zero mixed products. For the third quotient, relative
cohomology of $(\RP^\infty,\RP^1)$ injects into ordinary cohomology in
all degrees at least two, with image $t^2\F_2[t]$, and the relative cup
product is compatible with this injection. Integrally the space is
$\operatorname{Th}(2\gamma)$; this real plane bundle is oriented and is
the underlying real bundle of the complex sign line. If $u$ is its Thom
class and $c\in H^2(\RP^\infty;\Z)=\Z/2$ is its Euler class, the Thom
identity gives $u^k=u c^{k-1}$. Each $c^{k-1}$ is the nonzero class in
degree $2k-2$, so the assertion follows from the Thom isomorphism.
This also explains the integral nonnilpotence behind the mod-two
counterexample of Theorem~\ref{A:six-nonnilpotent}.
\end{proof}
